% AI-written, AI cross-reviewed; no human sign-off; not a claim that RH is proved
% RH Notes blueprint source excerpts for N318
% Verbatim excerpts in the notation of the parent blueprint. The labelled editions and source ranges appear below.
% See the unit YAML for the current proof status. Earlier-round passages retain their original wording.

% Start of atlas edition annotation. The following block is outside the original excerpt.
\paragraph{Status of named inputs in this reading.}
\noindent\nolinkurl{GlobalTheory.Published.HoweDualityReal} --- hypothesis at this head (\path{Development/AbelianVarietyTheory/PublishedTheorems/Divisors₂/HoweDualityReal.lean:21--31}); not compared.\par
\noindent\nolinkurl{GlobalTheory.Published.HeckeFunctionalEquationAbs} --- hypothesis at this head (\path{Development/AbelianVarietyTheory/PublishedTheorems/Divisors₂/HeckeFunctionalEquationAbs.lean:10--15}); not compared.\par
\noindent\nolinkurl{GlobalTheory.T.triEx_of_howe_duality_real} --- in-repo constant (\path{Development/AutomorphicArithmetic/ThetaCorrespondence/ThetaPairLifts/LocalTypeDictionary/TripleExistenceOfTheta.lean:49--50}); not compared.\par
\noindent\nolinkurl{GlobalTheory.FunPr.ThetaLocalInputs.howeDualityReal} --- record field (\path{Development/AutomorphicArithmetic/ThetaCorrespondence/ThetaPairLifts/LocalTypeDictionary/BernsteinUnitaryProductIrreducible.lean:57}); not compared.\par
% End of atlas edition annotation.

\section*{N318. Joined estimates at the square top term}

% Parent source SHA-256: 44e3b079775c8d7cc0d43b0314bf43a3959be56b2bbafa7f26e5441a1a9163de
% blueprint/ch_T.tex lines 376 through 404; r6:N318
\subsection{N318. A common bottom norm and a nonzero small coordinate}\label{r6:N318}
\textit{Partial.} Section~\ref{sec:T-r5-N318} gives the nine hypotheses, twelve conclusions, and the height and valuation bounds. Two bottom assertions have different tuple quantifiers. The norm formula holds at every eligible tuple. The simultaneous coordinate nonvanishing required by \nolinkurl{BotNe} concerns every bottom pair at the base tuple $\mu_0$.\footnote{\nolinkurl{normBig_alt} in \nolinkurl{Development/AutomorphicArithmetic/ThetaCorrespondence/ThetaPairLifts/LocalTypeDictionary/NormBigOfNonvanishing.lean}\allowbreak:14--41. \nolinkurl{bottomTheta_ne_zero_of_innerProduct_formula} in \nolinkurl{Development/AutomorphicArithmetic/ThetaCorrespondence/ThetaPairLifts/LocalTypeDictionary/BottomThetaNeZeroOfInnerProductFormula.lean}\allowbreak:19--43.}

\paragraph{The scalar in the norm formula.}
Put $a=P.a$, $d=P.d$ and $t_*=d-a>0$. For the Fock vector indexed by a minor, the anisotropic doubling calculation uses the constant automorphic function $1$, of trivial Satake parameter. On the right half-plane write $I(S)$ for the integral over the open subgroup unramified off $S$. Adding disjoint good places $U$ gives $I(S\cup U)=I(S)\prod_{v\in U}L_v(s+1/2)/d_v(s)$. The sets $S\cup U$ are cofinal, and the partial integrals tend to the full integral by integrability. The two convergent products have a nonzero normalizing denominator. Uniqueness of the limit gives the normalized Euler product times $I(S)$. At the bottom evaluation the anisotropic global formula and holomorphic values give the product of the two quotient volumes, the integral of the Weil coefficient on the open subgroup and
\[
 \frac{\prod_{k=0}^{a-1}\Lambda(t_*/2+a/2-k)}
 {d^S(t_*/2)}.
\]
Here the dictionary identifies the requisite partial Hecke values for the given $\Lambda$. Identify the volumes with masses. Sum the Fock terms, insert the block-level index and its local normalizing identity. Since $d=2h_0$ and $a=h_0+e_0$, the numerator runs through $J=[1-e_0,h_0]\cap\mathbb Z$ in reverse order. This gives $P.\mathrm{norm}=F(D,\delta,t_0,a,d,P.n)\prod_{s\in J}\Lambda(s)$ with the source's \nolinkurl{normFactor} $F$.\footnote{\nolinkurl{limit_eq_prod_mul}, \nolinkurl{eulerRight} in \nolinkurl{Development/AbelianVarietyTheory/Proofs/DecoratedLatticeProofs/DoublingEulerProductFarRight.lean}\allowbreak:24--43,75--91. \nolinkurl{rallis_inner_product_bottom_norm_eq} in \nolinkurl{Development/AutomorphicArithmetic/ThetaCorrespondence/ThetaPairLifts/LocalTypeDictionary/RallisInnerProductBottomNormEq.lean}\allowbreak:26--63. \nolinkurl{theta_lift_petersson_norm_eq_rallis_integral}, \nolinkurl{petersson_norm_bottom_theta_lift_eq} in \nolinkurl{Development/AutomorphicArithmetic/ThetaCorrespondence/ThetaPairLifts/LocalTypeDictionary/ThetaLiftPeterssonNormEqRallisIntegral.lean}\allowbreak:38--73,74--113.}

Fix $\delta,t_0,i,j$ before choosing any tuple. The ranks and the weight $P.n$ agree for every eligible bottom pair at these indices, so the $F$ factors agree. The factors are nonzero: the positive normalization, two nonzero masses, Fock sum and vacuum norm occur in the numerator, and the incomplete normalizer is nonzero at $\operatorname{Re}(t_*/2)>0$. Each $\Lambda(s)$ is nonzero by \nolinkurl{LNe} and $h_0\le\mathrm{halfRank}$. Intersect the eventual thresholds by maxima. If there is no eligible pair take the common scalar to be $1$. Otherwise choose one pair $P_*$ and one certified $\Lambda_*$ and take $A=P_*.\mathrm{norm}/\prod_{s\in J}\Lambda_*(s)$. The equality of the two $F$ factors gives the cross-product identity, so this $A\ne0$ works for every eligible tuple, pair, and certified $\Lambda$.\footnote{\nolinkurl{big_theta_petersson_norm_ratio_eq} in \nolinkurl{Development/AutomorphicArithmetic/ThetaCorrespondence/ThetaPairLifts/LocalTypeDictionary/ThetaLiftPeterssonNormEqRallisIntegral.lean}\allowbreak:115--131. \nolinkurl{big_theta_petersson_norm_ne_zero} in \nolinkurl{Development/AutomorphicArithmetic/ThetaCorrespondence/ThetaPairLifts/LocalTypeDictionary/BigThetaLiftNeZeroOfPeterssonNorm.lean}\allowbreak:14--38. \nolinkurl{normBig_alt} in \nolinkurl{Development/AutomorphicArithmetic/ThetaCorrespondence/ThetaPairLifts/LocalTypeDictionary/NormBigOfNonvanishing.lean}\allowbreak:14--41.}

\paragraph{Using the same minor on both sides.}
At $\mu_0$ the bottom modulus equals the prescribed conductor. The nonzero norm gives $h,T$ with $P.\mathrm{coord}(h,T)\ne0$. Choose one standard character of sign $-1$ and its Weil datum. The independent triangle input \nolinkurl{TriEx} supplies an initial triangle and its constant-function certificate. A change of basis gives the common triangle $X$ with $\beta=1$ and its transported certificate. Lift $h$ to the same finite adelic point $g$ for both formulas, using equality of the ambient ranks. Suppress the character, minor maps and $g$ in the notation. The class sums with inverse stabilizer orders give
\[
 \overline{P.\mathrm{coord}(h,T)}=\Omega u_A I_T,\qquad
 I'_{T'}=\Omega' u_B P'.\mathrm{coord}(h,T').
\]
In the first formula the class-sum factor times the uniform kernel factor is inverted, as is its tuple-specific away-unit. In the second formula the two uniform factors multiply. The real stabilizer weights commute with conjugation in the first sum.\footnote{\nolinkurl{bottomTheta_ne_zero_of_innerProduct_formula} in \nolinkurl{Development/AutomorphicArithmetic/ThetaCorrespondence/ThetaPairLifts/LocalTypeDictionary/BottomThetaNeZeroOfInnerProductFormula.lean}\allowbreak:29--62. \nolinkurl{rallis_bottom_formula_large_rank} in \nolinkurl{Development/AutomorphicArithmetic/ThetaCorrespondence/ThetaPairLifts/LocalTypeDictionary/RallisBottomFormulaLargeRank.lean}\allowbreak:20--43,54--66. \nolinkurl{rallis_bottom_formula_small_rank} in \nolinkurl{Development/AutomorphicArithmetic/ThetaCorrespondence/ThetaPairLifts/LocalTypeDictionary/RallisBottomFormulaSmallRank.lean}\allowbreak:18--28,35--51.}

For the matching, write $q,q'$ for the surjective polynomial maps to the same real-place model over $\mathbb C$. The small integral vanishes on $\ker q'$ and therefore factors as $J_s q'$ at tuple $s$. Before normalization the large integral is $\lambda_s J_s q$. The nonzero test integral at the distinguished polynomial $P'_0$ in that certificate rules out the identically zero branch, and the converse kernel inclusion makes $J_s$ injective. For the distinguished $P_0,P'_0$ the triangle relation gives $\lambda_s q(P_0)=c_0(s)q'(P'_0)$. At $s_1$ with $c_0(s_1)\ne0$ this forces $\lambda_{s_1}\ne0$. Hence $\kappa=c_0(s_1)/\lambda_{s_1}$ is fixed, and $q'(P'_0)\ne0$ yields $\lambda_s=c_0(s)/\kappa$ for every $s$. The signed complement of the minor indices gives an integral signed permutation $M$, and the common real-place factor $\rho$ gives
\[
 I_T=(\rho/\kappa)c_0(\mu_0)\sum_{T'}M_{TT'}I'_{T'} .
\]
If every small coordinate at the chosen $h$ vanished, this sum, and then the selected large coordinate, would vanish. This proves the two existential coordinate assertions even though the minor indices on the two sides can differ.\footnote{\nolinkurl{theta_bottom_minor_match} in \nolinkurl{Development/AutomorphicArithmetic/ThetaCorrespondence/ThetaPairLifts/LocalTypeDictionary/RallisBottomMinorMatching.lean}\allowbreak:27--62,145--189,190--203. \nolinkurl{bottomTheta_ne_zero_of_innerProduct_formula} in \nolinkurl{Development/AutomorphicArithmetic/ThetaCorrespondence/ThetaPairLifts/LocalTypeDictionary/BottomThetaNeZeroOfInnerProductFormula.lean}\allowbreak:59--71.}

\paragraph{Outstanding propositions.}
\textit{[Proved in Lean, not reproduced.]} The analytic input still to check is the anisotropic global Fock identity at the evaluation point, obtained from theta-tower compatibility and the anisotropic Siegel--Weil input, together with the split and inert local integral evaluations needed in the finite-place product above. On the small-coordinate side a specific remaining matching requirement is that, for each eligible $X,s$ and every compatible real/finite tensor factorization of its matching maps, the fixed finite vectors satisfy $r(X.f(s))=c\,r'(X.f'(s))$ for some scalar $c$. The construction passes from one reference matching datum to every matching datum via an intertwining isomorphism and its nonzero scalar. The bottom kernel-to-class-sum inputs for the two displayed coordinate formulas also retain their stated normalizations. \textit{[Hypothesis, discharged at the application.]} The nine outer premises retain the status recorded in Section~\ref{sec:T-r5-N318}; the smaller propositions here occur inside their conditional consequence.\footnote{\nolinkurl{rallis_inner_product_bottom_norm_eq} in \nolinkurl{Development/AutomorphicArithmetic/ThetaCorrespondence/ThetaPairLifts/LocalTypeDictionary/RallisInnerProductBottomNormEq.lean}\allowbreak:40--53. \nolinkurl{theta_bottom_finite_match} in \nolinkurl{Development/AutomorphicArithmetic/ThetaCorrespondence/ThetaPairLifts/LocalTypeDictionary/ThetaBottomFinitePlaceMatching.lean}\allowbreak:15--37. \nolinkurl{exists_matching_finiteData_bottomLevel} in \nolinkurl{Development/AutomorphicArithmetic/ThetaCorrespondence/ThetaPairLifts/LocalTypeDictionary/ExistsMatchingFiniteDataBottomLevel.lean}\allowbreak:16--24. \nolinkurl{rallis_bottom_formula_large_rank} in \nolinkurl{Development/AutomorphicArithmetic/ThetaCorrespondence/ThetaPairLifts/LocalTypeDictionary/RallisBottomFormulaLargeRank.lean}\allowbreak:28--43. \nolinkurl{rallis_bottom_formula_small_rank} in \nolinkurl{Development/AutomorphicArithmetic/ThetaCorrespondence/ThetaPairLifts/LocalTypeDictionary/RallisBottomFormulaSmallRank.lean}\allowbreak:26--39. \nolinkurl{square_top_term_joined_estimates} in \nolinkurl{Development/AutomorphicArithmetic/ThetaCorrespondence/ThetaPairLifts/LocalTypeDictionary/SquareTopTermJoinedEstimates.lean}\allowbreak:13--35.}

